\documentclass[11pt]{article}
\usepackage{fontspec,amsmath,amssymb,amsthm,hyperref}
\usepackage[margin=1in]{geometry}
\setlength{\emergencystretch}{3em}
\title{Normalized increments are independent standard Gaussians}
\author{LangevinLean mathematical exposition}
\date{}
\begin{document}\maketitle
\noindent Source module: \texttt{UnadjustedLangevin/BrownianInnovations.lean}. Compile this proof document with XeLaTeX or LuaLaTeX. The accompanying HTML pairs the mathematical argument with the exact Lean statements and records the scope of the formalization.


For a pre-Brownian coordinate and positive grid spacing \(h\), \(Z_n=(B_{(n+1)h}-B_{nh})/\sqrt h\) has standard normal law. Across independent coordinate paths, the entire coordinate-time family is independent.


\section{normalizedBrownianIncrement}


Define \(\xi_n=(B_{(n+1)h}-B_{nh})/\sqrt h\). The denominator is interpreted as real scalar division; positive-step law results below explicitly require \(h>0\), so that the variance rescaling is legitimate.



\noindent Lean declarations: \texttt{normalizedBrownianIncrement}.

\section{normalizedBrownianIncrement hasLaw}


A Brownian increment over an interval of length \(h\) is \(\mathcal N(0,h)\). Scaling by \(1/\sqrt h\) multiplies its variance by \(1/h\), producing \(\mathcal N(0,1)\). The proof rewrites \((n+1)h-nh=h\) and uses \((\sqrt h)^2=h\), justified by \(h>0\).



\noindent Lean declarations: \texttt{normalizedBrownianIncrement\_hasLaw}.

\section{normalizedBrownianIncrement memLp}


The exact standard Gaussian law provides every finite moment, in particular the second. Pull this square integrability back through \texttt{normalizedBrownianIncrement\_hasLaw} to prove \(\xi_n\in L^2(P)\).



\noindent Lean declarations: \texttt{normalizedBrownianIncrement\_memLp}.

\section{normalizedBrownianIncrement second moment}


Transfer the integral of \(x^2\) along the standard Gaussian law of \(\xi_n\). Its mean is zero and variance one, hence \(\mathbb E\xi_n^2=1\). This is an exact second-moment identity for the normalized increment.



\noindent Lean declarations: \texttt{normalizedBrownianIncrement\_second\_moment}.

\section{normalizedBrownianIncrement independent}


Brownian increments on the consecutive grid intervals are independent. Normalizing each by the same deterministic scalar is a coordinatewise measurable transformation and preserves independence. The formulation permits \(h=0\), where every increment is the constant zero and the independence assertion remains valid.



\noindent Lean declarations: \texttt{normalizedBrownianIncrement\_independent}.

\section{iIndepFun uncurry nat of aemeasurable}


Assume the row-valued random maps \((X_{i,n})_n\) are independent and each row is internally independent. Replace the almost-everywhere measurable coordinates by measurable versions; countability of \(\mathbb N\) lets all replacements in a row hold on one full-measure set. Apply independence of independent families, then transfer the result back across the coordinatewise almost-everywhere equalities.



\noindent Lean declarations: \texttt{iIndepFun\_uncurry\_nat\_of\_aemeasurable}.

\section{normalizedBrownianIncrements independent}


Each grid increment depends measurably on two entries of its coordinate grid path. Independence of those paths gives independence across coordinate rows, while Brownian independent increments gives independence within each row. Apply \texttt{iIndepFun\_uncurry\_nat\_of\_aemeasurable} to obtain joint independence of all \((i,n)\) innovations.



\noindent Lean declarations: \texttt{normalizedBrownianIncrements\_independent}.

\section{normalizedBrownianIncrements independent of paths}


Restrict each independent full coordinate path to the grid \(nh\), preserving independence under composition. The preceding theorem then yields joint coordinate-time independence of the normalized increments. This avoids requiring a separate measurable map into an uncountable path space.



\noindent Lean declarations: \texttt{normalizedBrownianIncrements\_independent\_of\_paths}.

\section{The entire innovation law}


Every coordinate-time innovation has law \(\mathcal N(0,1)\), and the family is jointly independent. The countable product-law theorem therefore identifies the whole innovation map with \(\bigotimes_{(i,n)}\mathcal N(0,1)\). Countability supplies measurability of the path-valued map from coordinatewise almost-everywhere measurability.



\noindent Lean declarations: \texttt{normalizedBrownianIncrements\_hasLaw}.

\end{document}
